\section{The two inner arrows in the theta tower}

% BEGIN I77A1
\paragraph{I77A1. The Fock projection and the local real contraction.}
Fix a real place of a step $k+1<m$. Write $H_1,H_2$ for the two
covering discrete series and $\mathcal F$ for the Hilbert Fock model
there. The signed infinitesimal parameters of the source and target
are the duals of the original two parameters, with signatures
\[
 (|\alpha|+|\beta|,|\gamma|+|\delta|),\qquad
 (|\alpha|+t_0+|\delta|,|\gamma|+|\beta|).
\]
The second positive block includes the prescribed string of length
$t_0$. The node gap gives the separated parameters to which the
source's archimedean results apply. Write $\iota$ for polynomial
realization in $\mathcal F$.

Two local results used below have distinct jobs. The first supplies a
joint highest vector $\Phi$ and a nonzero projected vector. More
precisely, for a unit highest vector $v$ of $H_1$, formal degree $d$
for the given Haar measure and $m_v(g)=\langle v,\pi_1(g)v\rangle$,
let $E=d\int\overline{m_v(g)}\,\omega(g,1)\,dg$. The local
coefficient-decay estimate makes this an integral of bounded operators
on each vector. Discrete-series coefficient orthogonality makes $E$
an orthogonal projection (convolution of its coefficient with itself
reproduces the coefficient). Hence
\[
 \int \langle\iota\Phi,\omega(g,1)\iota\Phi\rangle
                  \overline{m_v(g)}\,dg
        =d^{-1}\langle\iota\Phi,E\iota\Phi\rangle
        =d^{-1}\|E\iota\Phi\|^2 .
\]
The strict inequality can be checked before constructing the second
discrete series in $\mathcal F$. Let $L$ be the common irreducible
compact type generated by the two unit highest vectors,
$e=\dim L$, and let $j_1,j_2$ embed $L$ in the Fock and discrete
models, respectively, carrying a fixed unit vector $v_0$ to their
chosen vectors. These equivariant embeddings can be made isometric
since the invariant inner product on $L$ is unique up to scalar.
Put $A_i(g)=j_i^*\pi_i(g)j_i$ with the Fock action in the first
slot. For $g=k\exp Y$ in the polar decomposition the source's
Cartan coefficient estimates make both $A_i(\exp Y)$ positive
semidefinite. Thus
$\operatorname{tr}(A_2(g)^* A_1(g))\ge0$: the common compact factor
cancels, and for positive $C,D$ the trace of $CD$ is the trace of
$D^{1/2}CD^{1/2}$. It is strictly positive at the identity and
therefore on a Haar-positive open set.

The Haar integral of the original product of two scalar coefficients
equals the integral of its left and right compact averages. Fubini
applies by integrability and Haar invariance. For fixed $g$, averaging
over $k_1,k_2$ replaces the two rank-one projectors on $k_1^{-1}v_0$
and $k_2v_0$ in that product by $1_L/e$. Indeed their averages
commute with the irreducible compact action and have trace one.
The resulting function is precisely
$e^{-2}\operatorname{tr}(A_2(g)^*A_1(g))$. Its integral is strictly
positive by the preceding paragraph, which proves $E\iota\Phi\ne0$.
This uses positivity on the two compressed Cartan coefficients,
proved in the source for the harmonic and discrete-series types.
The second local result supplies an equivariant bilinear isometry
$T:H_1\times H_2\to\mathcal F$ on pure tensors, with
$\langle T(x,y),T(x',y')\rangle=
\langle x,x'\rangle\langle y,y'\rangle$.
These are the source's archimedean projection and embedding theorems
for the indicated parameters.

For fixed $y$ put $T_y x=T(x,y)$ and define
$B(P,y)=T_y^*\iota P$. This is linear in $P$, conjugate-linear in
$y$, and $\|B(P,y)\|\le\|\iota P\|\,\|y\|$. If $y\ne0$ and all these
values vanished, $T_y x$ would be orthogonal to the dense polynomial
subspace for every $x$, contradicting
$\|T_y x\|=\|x\|\,\|y\|$ at any $x\ne0$. Equivariance and
unitarity give
\[
 \pi_1(g)T_y^*=T_y^*\omega(g,1),\qquad
 T_{\pi_2(h)y}^*=T_y^*\omega(1,h^{-1}).
\]
Since the maximal compact subgroup preserves polynomial degree, the
orbit of $B(P,y)$ in the first group has finite-dimensional span;
thus it lies in the Harish--Chandra module, hence in the smooth model.

Let $U$ be the source linear isometry with $H_1$, $U'$ the target
\emph{antilinear} isometry with $H_2$, and $\kappa$ the isomorphism
of the two polynomial coordinate systems. The actual local map is
\[
             A(P,y)=U^{-1}B(\kappa P,U'y).
\]
The two conjugations make this map complex bilinear. A nonzero smooth
input from the target Harish--Chandra module and the density argument
above supply a nonzero output for some $P$. The two group identities
give the infinitesimal rules by differentiation; the source model is
transported by its conjugated matrix and the target by its plain
matrix. The derivatives of their covering characters contribute the
scalar trace terms $-a_T\operatorname{tr}X$ and
$+a_S\operatorname{tr}Y$, respectively, where $m_1/2=a_T$ and
$m_2/2=a_S$. This fixes the contraction once the local models, their isometries
and the embedding have been chosen, before any twist is specified.
% END I77A1

% BEGIN I77A2
\paragraph{I77A2. Tensoring the contractions and testing the full group.}
For a product $P=\prod_\tau P_\tau$ in the disjoint polynomial
variables and a pure tensor $u=\otimes_\tau u_\tau$ set
\[
                 I(P,u)=\bigotimes_\tau A_\tau(P_\tau,u_\tau).
\]
Expansion in monomials identifies the global polynomial algebra
with the tensor of the local polynomial algebras, so this extends
unambiguously to all $P$ and all algebraic tensors of smooth vectors.
The local nonzero tests in I77A1 give a nonzero pure output, witnessed
by a product of linear functionals nonzero on all its factors.
For $P=\sum_j\prod_\tau P_{j,\tau}$ and a pure Hilbert test
$w=\otimes w_\tau$, a bound for its matrix coefficient is
\[
 |\langle w,I(P,u)\rangle|
 \le \left(\sum_j\prod_\tau
             \|w_\tau\|\,\|\iota_\tau\kappa_\tau P_{j,\tau}\|\right)
                \prod_\tau\|u_\tau\|.
\]
The local derivative rules imply the global rules first on pure
tensors. For the target rule on a general tensor, the scalar
functional $u\mapsto\langle w,I(P,u)\rangle$ is a finite sum of
products of bounded local linear functionals. By Riesz each product
is a pure Hilbert test on the target; hence the weak derivative in
the definition can be applied to this functional. Injectivity of
the realized polynomial map identifies a proposed pointwise
derivative with the corresponding source or target polynomial
operator by uniqueness of scalar derivatives.

There is a further point because the equivariance conditions test
\emph{all} real unitary group elements that send one realized
polynomial to another. Here is the finite-degree argument. In the
unitary frame let $J=\operatorname{diag}(1_p,-1_q)$, let $g^*Jg=J$,
and suppose $g$ carries a nonzero realized $P$ to a realized $P'$.
The conjugate of the degree-$\le\deg P'$ subspace under $g^{-1}$
is finite dimensional, contains the realized $P$, and is invariant
under $\exp(sY)$ with $Y=g^{-1}(\mathrm iJ)g=\mathrm iJg^*g$.
Indeed this one-parameter group is conjugate to a compact subgroup
preserving degree. Its infinitesimal polynomial operator therefore
has all its powers on $P$ in a finite-dimensional space: differentiate
the pointwise orbit inside the finite-dimensional space, which is
closed in the Hausdorff pointwise topology, and use injectivity of
realization. If $g^*g$ commuted with $J$, then $(g^*g)J(g^*g)=J$
would imply $(g^*g)^2=1$. Positivity would give $g^*g=1$ and $g$
would already be in the maximal compact subgroup.
Otherwise $Y$ has a nonzero pair of opposite off-block entries.
The Fock Lie formula in either member has degree-raising part
multiplication by a nonzero quadratic polynomial $Q_Y$ in this case
(the other member has positive rank). All other terms preserve or
lower the degree. The leading homogeneous term of the $j$th power
on $P$ is $Q_Y^j P_{\rm top}$ and has degree $\deg P+2j$.
The polynomial domain makes all these terms nonzero, a contradiction.

It therefore suffices to integrate the two infinitesimal identities
on the maximal compact group. Each block-unitary matrix is an
exponential there by diagonalizing its unitary blocks. Along such
a one-parameter group the Fock polynomial orbit remains in a fixed
degree space. The covariance identity follows by differentiating the
matrix coefficient after moving the inverse action to the other side;
the two derivative terms cancel and its value at zero is the
identity. Pure Hilbert tests separate algebraic tensors. For the
zero polynomial case injectivity gives $P'=0$. This proves the two
full group tests, as well as the boundedness, two Lie tests and
nonvanishing required of the original real intertwiner. For
$k+1<m$ the model indices $\min(k,m),\min(k+1,m)$ are $k,k+1$.
The local construction and its tensor use only these fixed models and
their character exponents. Every twist's vacuum realizes the same
polynomial Lie and compact laws, so this single operator witnesses
$R.\mathrm{ex\_ItY}$ under the local inputs above.
% END I77A2

% BEGIN I77A3
\paragraph{I77A3. Averaging the stored intertwiner together with the integral.}
Fix $k+1<m$ and a twist, and write $I$ for the stored operator certified
by I77A2 and weak choice. Write $B_F(P,y)$ for the theta integral
of a fixed input multiplicity form $F$ as a function of the output
group variable. Assume the absolute integrability of its compact
averages as in the theta kernel estimates. At a fixed real place take
normalized Haar measure on the target maximal compact group $K'$.
In the finite-dimensional compact orbit span of a polynomial write
$kP=\sum_a c_a(k)P_a$. The $c_a$ are matrix coefficients of one
finite-dimensional compact representation. For a smooth local input
$y$ set $y_a=\int_{K'} c_a(k)\pi'(k)y\,dk$.
The translates of each $c_a$ span a finite-dimensional space; changing
variables shows the same for the translates of $y_a$, so $y_a$ is
$K'$-finite and is a smooth vector in the given model. Invariance
of the simultaneous target action gives
\[
       (I(P,y),B_F(P,y))=\sum_a (I(P_a,y_a),B_F(P_a,y_a)).
\]
All other places are held fixed in this identity. For $I$ the
average can be passed through each Hilbert coefficient by its boundedness
in the varied input. For $B_F$ one uses the source's absolutely
integrable compact averages of the theta kernel times a cusp form.
Translate the integration variable and apply Fubini to each of the
finitely many coefficient functions. Iterating through the finite
set of real places and using linearity on algebraic tensors proves
that every simultaneous pair belongs to the span of its values on
$K$-finite inputs. Since $I$ is nonzero,
its restriction to finite vectors is therefore nonzero as well.

% END I77A3

% BEGIN I77A4
\paragraph{I77A4. Kernel inclusion for the particular stored intertwiner.}
Fix $(k,t)$, where $k+1<m$, and write $I$ for the intertwiner actually
stored at $k$. I77A2 and the weak choice make it nonzero and give all
its covariance identities. For an input multiplicity form $F$, denote
the desired integral on a polynomial $P$ and a target smooth vector
$y$ by $B_F(P,y)$. The difficulty is that linear relations among
values of $I$ must also hold among these integrals. On the
$K$-finite subspaces $Y_j$ of the target and $V_j$ of the source
at the real places $j$, one can prove this simultaneously for every
relation.

At one place the source uses the downward archimedean correspondence
for these exact parameters. Its relevant assertion is that a nonzero
polynomial intertwiner into $A\otimes Y_j$, for a unitarizable
Harish--Chandra module $A$, has image $N\otimes Y_j$ with
$N\subset A$ isomorphic to the irreducible source module $V_j$.
Here is how it provides the local kernel test. Let $b_1,b_2$ be two
local bilinear pairings into inner-product spaces satisfying the
source Lie skew-adjointness identity and target Lie invariance.
Suppose $b_1\ne0$. On the covering models the second variable is
conjugate-linear. Take the pair $(b_1,b_2)$ and the span of its
values, with the direct-sum inner product. The skew identity defines
the source Lie action there: applying a generator to a relation
gives a vector orthogonal to the spanning vectors, hence zero.
The source's admissibility and finite-support result realizes the
target-invariant pair as contraction of a polynomial intertwiner
into that unitary module tensored with $Y_j$. This finite-support
step is needed for the algebraic tensor product. The downward
assertion now puts every value in one irreducible $N$. The kernel
of projection to the first component inside $N$ is stable: the
skew identity tested against the image of its own differentiated
vector makes the squared norm of that image zero. The projection
is nonzero by $b_1\ne0$, hence injective on $N$. In particular,
\[
       \sum_a b_1(P_a,y_a)=0\quad\Longrightarrow\quad
                      \sum_a b_2(P_a,y_a)=0 .
\]
Transfer through the conjugate target isometry gives this test in
the actual models.

Apply the test to local slices obtained by holding all other
polynomials and vectors fixed. Write $V=\bigotimes_j V_j$.
Choose a nonzero pure-input value of $I$ on finite vectors. Its
existence follows from the simultaneous compact averaging proved
in I77A3. For each $j$ contract the other factors of that
value by linear functionals to give a nonzero element of $V_j$.
Such contractions separate tensors, by expanding any given algebraic
tensor in finite bases. Denote the resulting local slice of $I$
by $A_j$. Its values span $V_j$: their span is nonzero and stable
under the source Lie algebra, since differentiation here acts on
the polynomial input. All contracted local slices of $I$ satisfy
the local test. For the slices of $B_F$, use the injection of smooth
cuspidal functions of uniform moderate growth into the quotient
$L^2$. Rapid decay makes the functions and the required derivatives
bounded on the finite-volume quotient. Translation invariance gives
\[
  \langle R_X f,h\rangle+\langle f,R_X h\rangle=0.
\]
Indeed differentiate the integral of
$\overline{f(g\exp(sX))}h(g\exp(sX))$; the two products of
bounded derivatives and functions give an integrable uniform bound.
A continuous rationally invariant function with zero $L^2$ class is
zero: its nonzero locus is open, and nullity on a fundamental domain
extends by rational invariance to Haar nullity on the group.
Thus the resulting kernel inclusion holds for the actual functions.
The target Lie identity follows from differentiating the theta
integral in its integration variable with both factors translated.

Let $q_j:\mathcal P_j\otimes Y_j\twoheadrightarrow V_j$ be the
linearization of $A_j$. The local tests say that $I$ and $B_F$ on
finite vectors vanish when a relation in $\ker q_j$ is inserted into
any one slot. For $I$ this follows after every contraction in the
other output factors, which suffices by separation. The quotient of
the tensor of the local input spaces by the sum of these kernels
in each slot is $\bigotimes_j V_j$; this follows immediately by
choosing a complement of each $\ker q_j$. Hence
$I=\Phi\circ(\bigotimes q_j)$ and $B_F=T_0\circ(\bigotimes q_j)$.
The endomorphism $\Phi$ commutes with each source Lie action.
To check it is scalar, hold every input factor except $j$ fixed
and contract all other output factors. The induced endomorphism
of $V_j$ commutes with its Lie action and is scalar by absolute
irreducibility. As contractions separate tensors, $\Phi$ commutes
with all endomorphisms in each individual factor. Products of
rank-one endomorphisms then show that it is scalar on the whole
tensor: project onto one pure line and then move that line to an
arbitrary pure line. Write $\Phi=c\,1$. The nonzero test for $I$
gives $c\ne0$. Consequently $I$ spans the full source on finite
vectors and $B_F=(c^{-1}T_0)I$ there. This proves factorization for
the stored normalization of the intertwiner.
% END I77A4

% BEGIN I77A5
\paragraph{I77A5. The smooth multiplicity arrow.}
For clarity about the analytic inputs, the theta integral in this step
is
\[
 B_F(P,y)(g)=\int_{\mathcal F_{k+1}}
      \theta(g,h;\mathrm{pt}_{t,k}(P))\,F(y)(h)\,d\nu_{t,k+1}(h).
\]
The chain's kernel estimates and cuspidality theorem supply absolute
integrability, smoothness, uniform moderate growth and cuspidality
of $B_F$ and the two Lie differentiation laws on the indicated
spaces. Arbitrary smooth pure inputs extend to algebraic tensors by
linearity and absolute integrability. Finite translations give
the required level character by moving the right translation of
the output variable to the finite vacuum on the same member of the pair. On the
finite vectors I77A4 produces a map $T$ with $T I=B_F$. Its values
are smooth, cuspidal and of uniform moderate growth, as every finite
source vector is a sum of $I(P,y)$. On those generators, compact
equivariance and Lie differentiation follow by transferring the
identities for $I,B_F$; so does the finite level law. They extend to
their span. Use the source's globalization of such a
$(\mathfrak g,K)$-map into cusp forms to get a smooth multiplicity
map $U_F$ on the prescribed source representation. This is the
Casselman--Wallach input at this step, with the moderate growth,
level and square-integrability hypotheses just described.

The graph of $U_F$ is a linear subspace. It contains
$(I(P,y),B_F(P,y))$ for all finite $y$ and hence, by the averaging
identity for the pair, for every smooth $y$ in the domain. No
continuity of an unspecified extension of $I$ in a completion is
needed here. To obtain a linear lift in $F$, choose a vector-space
basis of the input multiplicity space, use these $U_F$ on the basis,
and extend linearly. The target is the required multiplicity
subspace, and bilinearity of the absolutely convergent integral gives
the required identity and integrability for every $F$. This
constructs $R.\mathrm{ex\_ioY}$ with the prescribed fundamental
domain and the intertwiner actually stored in $R$.

Thus the two inner entries in E77A6 reduce on the original data to
the positive archimedean projection and tensor embedding, the
downward local image theorem and finite-support realization, the
theta cusp and regularity estimates and the specified globalization,
together with the model and vacuum identities. The formal source
provides these from its published inputs. The constructions in
I77A1--I77A5 supply the contraction, its group tests, factorization
and the passage to smooth vectors. The cited archimedean correspondence
and global theta estimates remain analytic inputs to this account.
% END I77A5
